\documentclass[11pt,letterpaper]{article}
\usepackage[utf8]{inputenc}

%----- Configuración del estilo del documento------%
\usepackage{epsfig,graphicx}
\usepackage[left=2cm,right=2cm,top=1.8cm,bottom=2.3cm]{geometry}
\usepackage{fancyhdr}
\usepackage{lastpage}
\usepackage{url}
\pagestyle{fancy}
\fancyhf{}
\rfoot{\textit{Página \thepage \hspace{1pt} de \pageref{LastPage}}}


%------ Paquetes matemáticos básicos --------%
\usepackage{amsmath}
\usepackage{amssymb}
\usepackage{amsthm}

\usepackage[spanish]{babel}
\usepackage{graphicx}
\usepackage{hyperref}
\usepackage{cancel}
\usepackage{tabularx}
\usepackage{xcolor}
\usepackage[table]{xcolor}
\usepackage{colortbl}
\usepackage{array, multirow, multicol, tabularx}
\usepackage{tcolorbox}
\newtheorem{theorem}{Theorem}[section]
\newtheorem{corollary}{Corollary}[theorem]
\newtheorem{lemma}[theorem]{Lemma}

%------si-------%
\definecolor{B}{HTML}{FFFFFF}
\definecolor{G}{HTML}{5e5e5e}
\definecolor{R2}{HTML}{d53d40}
\definecolor{A2}{HTML}{034190}
\definecolor{V2}{HTML}{7faa50}
\newcommand{\R}{\mathbb{R}}
\newcommand{\C}{\mathcal{C}}
\newcommand{\N}{\mathbb{N}}
\newcommand{\Z}{\mathbb{Z}}
\newcommand{\Q}{\mathbb{Q}}
\renewcommand{\theenumi}{\Roman{enumi}}
\renewcommand{\labelenumi}{{\theenumi}.}

\begin{document}

%------ Encabezado -------- %
\begin{center}
    \begin{minipage}{3cm}
    	\begin{center}
    		\includegraphics[height=3.4cm]{logo_unam.png}
    	\end{center}
    \end{minipage}\hfill
    \begin{minipage}{10cm}
    	\begin{center}
    	\textbf{\large Universidad Nacional Autónoma de México}\\[0.1cm]
        \textbf{Facultad de Ciencias}\\[0.1cm]
        \textbf{Selectivo OMF 2015}\\[0.1cm]
        Problema 2 (Oscilaciones)\\[0.1cm]
         El\'ias L\'opez Rivera\\[0.1cm]
        \texttt{ elias.lopezr\,@ciencias.unam.mx }\\[0.1cm]
        Fecha:\,\,31/07/2025
    	\end{center}
    \end{minipage}\hfill
    \begin{minipage}{3cm}
    	\begin{center}
    		\includegraphics[height=3.4cm]{Logo_FC.png}
    	\end{center}
    \end{minipage}
\end{center}

\rule{17cm}{0.1mm}

%------ Fin de encabezado -------- %
\begin{tcolorbox}[
	title = \textcolor{black}{\textcolor{white}{Problema 2}}]
\textit{El sistema descrito en la figura de abajo consta de dos esferas de masa $m$, cada una de ellas esta sujeta a un par de barras que pueden girar libremente
haciendo un \'angulo $\theta$ con la vertical. Anclado a la par de barras inferiores se encuentra un bloque de masa $M$. El dispositivo gira con una frecuencia \'angular 
$\omega$ alrededor del eje vertical como se muestra en la figura.El sistema completo esta sujeto a la fuerza gravitacional.\,\\
\,\\
\begin{center}
    \includegraphics[height=0.4\textwidth]{Captura de pantalla (214).png}
\end{center}\,\\
\,\\
Desde un sistema que gira con las barras, haz un diagrama que actu\'an sobre cada una 
de las masas y escribe las ecuaciones de movimiento de cada una de ellas.\,\\
\,\\
\textbf{a)} Encuentra el \'angulo $\theta_0$ en el cual el sistema esta en equilibrio\,\\
\,\\
\textbf{b)} Si el sistema completo se desplaza un pequeño \'angulo alrededor del equilibrio $\theta_0$ oscilar\'a con cierta frecuencia
$\Omega$. Determina la frecuencia de oscilaci\'on $\Omega$ 
}
\end{tcolorbox}
\newpage
\,\\
\textbf{Soluci\'on}
\begin{figure}
     \centering
    \includegraphics[height=0.5\textwidth]{IMG_1314.jpg}
    \caption{Diagrama de cuerpo libre de la masa $m$}
    \label{Cono rebanado}
\end{figure}\,\\
\,\\
Si nos posicionamos sobre un sistema de referencia que rota con velocidad \'angular $\omega$, tendremos que el movimiento sera solo en dos dimensiones, proponemos un sistema de referencia en coordenadas polares
en donde el polo sera la posici\'on de la masa $m$, la primera relaci\'on importante que notamos es que $\alpha=90-\theta$, diferenciando obtenemos $\ddot{\alpha}=\ddot{\theta}$, finalmente del diagrama de fuerzas obtenemos
\,\\
\,\\
\begin{align*}
    F_{\alpha}&=m(l \ddot{\alpha}-\cancelto{0}{2\dot{l}\dot{\alpha}})=mg\sen(\theta)+T_2cos(2\theta-90)-m\omega_0^2l\sen(\theta)cos(\theta)\\
    \,\\
    &=mg\sen(\theta)+T_2\sen(2\theta)-m\omega_0^2l\sen(\theta)cos(\theta)
\end{align*}\,\\
Aplicando nuestra relaci\'on entre las tasas de cambio respecto al tiempo de $\alpha$ y $\theta$ se tiene que:\,\\
\,\\
\begin{equation*}
    2ml\ddot{\theta}=m\omega_0^2l\sen(2\theta)-2mg\sen(\theta)-2T_2\sen(2\theta)
\end{equation*}\,\\
Posteriormente hacemos un digrama de cuerpo libre de la masa $M$ en coordenadas polares, cuyo polo sera la posici\'on de la 
masa $M$, de donde obtenemos las siguientes ecuaciones:
\begin{figure}
     \centering
    \includegraphics[height=0.5\textwidth]{IMG_1316.jpg}
    \caption{Diagrama de cuerpo libre de la masa $M$}
    \label{Cono rebanado}
\end{figure}\,\\
\,\\
\begin{align*}
    \vec{r_M}=2l\cos(\theta)\implies \ddot{r}&=-2l(\dot{\theta}^2\cos(\theta)+\ddot{\theta}\sen(\theta))\,\\
    \,\\
    F_{r_M}=M(\ddot{r_M}-\cancelto{0}{r_M\dot{\alpha}^2}\,\,\,\,)&=-2ML(\dot{\theta}^2\cos(\theta)+\ddot{\theta}\sen(\theta))\,\\
    \,\\
    Mg-2T_2\cos(\theta)&=-2Ml(\dot{\theta}^2\cos(\theta)+\ddot{\theta}\sen(\theta))\,\\
    \,\\
    2Mg\sen(\theta)-2T_2\sen(2\theta)&=4Ml\ddot{\theta}\sen^2(\theta)+2Ml\dot{\theta}\sen(2\theta)
\end{align*}\,\\
Finalmente combinando ambas ecuaciones de movimiento obtenemos que:\,\\
\begin{align*}
    2(m+2M\sen^2(\theta))l\ddot{\theta}+2Ml\dot{\theta}^2\sen(2\theta)&-ml\omega_0^2\sen(2\theta)+2(M+m)gsen(\theta)\,\\
    \,\\
    &=2T_2\sen(2\theta)-2T_2\sen(2\theta)=0
\end{align*}\,\\
Aplicando la condici\'on de equilibrio $\dot{\theta}=\ddot{\theta}=0$, obtenemos la siguiente ecuaci\'on para $\theta_0$:\,\\
\,\\
\begin{align*}
    2(M+m)g\sen(\theta_0)=ml\omega_0^2\sen(2\theta)_0\implies\cos(\theta_0)=\frac{2g(M+m)}{ml\omega_0^2}
\end{align*}\,\\
Finalmente analicemos la ecuaci\'on de movimiento para pequeñas oscilaciones  alrededor de $\theta_0$, tenemos que $\theta=\theta_0+\epsilon$ ($\dot{\theta}=\dot{\epsilon},\ddot{\theta}=\ddot{\epsilon}$), con $\epsilon$ muy pequeño, luego aproximando hasta potencias de primer orden:\,\\
\,\\
\begin{align*}
    \sen(\theta_0+\epsilon)&=\sen(\theta_0)\cos(\epsilon)+\cos(\theta_0)\sen(\epsilon)\approx \sen(\theta_0)+\epsilon\cos(\theta_0)\,\\
    \,\\
    \sen^2(\theta_0+\epsilon)&\approx \sen^2(\theta_0)+\epsilon\sen(2\theta_0) 
\end{align*}\,\\
Finalmente valuando en la expresi\'on principal obtenemos:\,\\
\,\\
\begin{align*}
    2(m+M(\sen^2(\theta_0)+\epsilon\sen(2\theta_0)))l\ddot{\epsilon}+4Ml\dot{\epsilon}(\sen(2\theta_0)+\epsilon\cos(2\theta_0))\\
    \,\\
    -ml\omega_0^2(\sen(2\theta_0)+\epsilon\cos(\theta_0))+2(M+m)g(\sen(\theta_0)+\epsilon\cos(\theta_0))=0
\end{align*}\,\\
Finalmente aproximando todo a potencias de primer orden en las varibales $\epsilon,\dot{\epsilon},\ddot{\epsilon}$, y recordando la definici\'on de $\theta_0$ tenemos que:
\,\\
\begin{equation*}
    2l(m+2M\sen^2(\theta_0))\ddot{\epsilon}-\epsilon(2ml\omega^2_0\cos(2\theta_0)-2(m+M)g\cos(\theta_0))=0
\end{equation*}\,\\
\,\\
De donde obtendremos que:\,\\
\,\\
\begin{equation*}
    \ddot{\epsilon}+\frac{2(m+M)g\cos(\theta_0)-2ml\omega^2_0\cos(2\theta_0)}{l(m+2M\sen^2(\theta_0))}\epsilon=0
\end{equation*}\,\\
Notemos que para que haya equilibrio estable y la ecuaci\'on describa un M.A.S, es necesario que:\,\\
\,\\
\begin{equation*}
    2(m+M)g\cos(\theta_0)-2ml\omega_0^2\cos(2\theta_0)>0
\end{equation*}\,\\
Finalmente obtenemos la frecuencia para este caso:\,\\
\,\\
\begin{equation*}
    \Omega=\frac{2(m+M)g\cos(\theta_0)-2ml\omega^2_0\cos(2\theta_0)}{l(m+2M\sen^2(\theta_0))}
\end{equation*}

\end{document}